\section{Verification purpose}

\textbf{Case ID:} \texttt{\CaseID}. This case verifies coupling among the simulation clock, solar burn-period calculation, and surface-fire diurnal multiplier. With uniform fuel, moisture, wind, and slope, the question is whether ELMFIRE reduces and restores radial advance at the prescribed clock hours, each within \SI{1}{hour}. The case exposes clock offsets, failed midnight wrapping, reversed interval tests, or an unapplied multiplier; it does not calibrate the multiplier.

\section{Mathematical formulation and numerical implementation}

For simulation time $t$ seconds and forecast start $H_0$, ELMFIRE uses
\begin{equation}H(t)=\operatorname{mod}\!\left(H_0+t/3600,24\right).\end{equation}
The solar routine calculates sunrise and sunset from the lower-left coordinate, year, and hour of year. With default center fraction $f_c=0.667$ and burn length $L_b=\SI{10}{hour}$,
\begin{align}H_c&=H_{\rm rise}+f_c(H_{\rm set}-H_{\rm rise}),\\H_b&=H_c-L_b/2,&H_e&=H_c+L_b/2.\end{align}
The stop hour is wrapped above 24. For this non-wrapping interval,
\begin{equation}D(t)=\begin{cases}1,&H_b\le H(t)\le H_e,\\0.1,&H(t)<H_b\ \text{or}\ H(t)>H_e.\end{cases}\end{equation}
The spread-rate calculation multiplies its unadjusted rate by $D(t)$. Uniform flat, wind-free conditions produce an approximately circular fire whose equivalent radius $r(t)=\sqrt{A(t)/\pi}$ has full and reduced slope plateaus.

\section{Derivation of expected behavior}

The controlled response specifies burn-period start \SI{7.5}{hour} and stop \SI{17.5}{hour}. Full-to-reduced and reduced-to-full transitions are therefore expected at those respective clock hours. Since $H_0=\SI{10}{hour}$, the events first occur approximately \SI{7.5}{hour} and \SI{21.5}{hour} after ignition and repeat during the three-day run. This oracle follows from the controlled targets and clock equations; calculated times are independently inferred from the arrival raster.

\section{Verification metrics and acceptance criteria}

The evaluator selects the lexically latest \texttt{time\_of\_arrival\_*} GeoTIFF, masks nodata and values outside $0\le t<10^{8}$ seconds, and requires 50 burned cells. Analysis ends at the earliest domain-edge arrival so boundary truncation cannot mimic slowdown. At \SI{30}{minute} intervals, burned-cell count is converted to equivalent radius and differenced. After radius exceeds three cells, rates below half the 80th percentile of positive rates are reduced. Median full-to-reduced and reduced-to-full clock hours are then calculated.

\begin{center}\begin{tabular}{@{}llll@{}}
\toprule Metric & Target & Acceptance rule & Role \\
\midrule
Reduction starts & \SI{17.5}{hour} & absolute error $\le\SI{1}{hour}$ & burn-period stop \\
Reduction ends & \SI{7.5}{hour} & absolute error $\le\SI{1}{hour}$ & burn-period start \\
\bottomrule\end{tabular}\end{center}

Both rows are required. A missing raster, insufficient burned area or duration, or undetected transition produces NOT EVALUABLE, never PASS.

\section{Simulation configuration and predicted outcome}

\begin{center}\begin{tabular}{@{}lll@{}}
\toprule Input & Value & Verification role \\
\midrule
Domain/grid & $126\times126$; \SI{30}{m} & circular-front resolution \\
CRS & EPSG:32635 & location for solar calculation \\
Fuel/terrain & uniform FM3; zero slope & constant baseline spread \\
Wind/dead moisture & zero; 6, 7, 8 percent & constant forcing \\
Live moisture & herb 60; woody 90 percent & constant forcing \\
Ignition & $(1890,1890)$ m at $t=0$ & centered radial growth \\
Time control & \SI{1}{s}; \SI{259200}{s} stop & three daily cycles \\
Calendar & year 2025; hour 2890; start 10 & clock coupling \\
Diurnal control & enabled; factor 0.1 & two rate plateaus \\
Output & final TOA, spread rate, FLIN & case-local evidence \\
\bottomrule\end{tabular}\end{center}

The predicted arrival field is concentric; contour spacing contracts overnight and expands during the burn period, with both transitions inside their acceptance bands.

\section{Actual simulation results}

The variant generated the required evidence. The evaluator selected \path{outputs/time_of_arrival_0000001_0266400.tif}, a $126\times126$, \SI{30}{m}, EPSG:32635 raster whose values are arrival seconds. The log reports eventual front stagnation on the finite domain followed by successful shutdown; the evaluator excludes edge contact and later stagnation.

Postprocessing evaluated \VariantCount{} variant and all \MetricCount{} required rows. The decision is \textbf{\VerificationStatus}, with \PassingMetricCount{} passing rows. Missing measurements would appear as NOT EVALUABLE.

\begin{figure}[H]
\centering
\IfFileExists{../figures/whole_domain_result.pdf}{\includegraphics[width=0.82\textwidth,height=0.58\textheight,keepaspectratio]{whole_domain_result.pdf}}{\fbox{\parbox{0.85\textwidth}{No case-local whole-domain arrival-time raster is available.}}}
\caption{Final time of arrival for variant \texttt{main}, in seconds, from \protect\path{outputs/time_of_arrival_0000001_0266400.tif}. The $3780\times3780$ m EPSG:32635 domain shows the near-circular geometry required by the radial metric. Acceptance is determined by the generated table.}
\end{figure}

\section{Calculated metrics and verification decision}

\begin{center}\begin{tabular}{@{}llll@{}}
\toprule Variant / metric & Expected & Calculated & Status \\
\midrule
\MetricRows
\bottomrule\end{tabular}\end{center}

This authoritative table and \textbf{\VerificationStatus} are generated by \path{scripts/postprocess.py} from \path{outputs/metrics.json}; the report does not recompute them. Comparison symbols remain in LaTeX math mode, so $\le$ and $\ge$ cannot be misparsed as table markup.

\section{Technical interpretation}

Both required transitions are evaluable and within their one-hour bands. This supports the forecast-hour offset, midnight wrapping, solar burn-period selection, and application of the factor to surface spread. Half-hour sampling limits timing resolution. A common transition shift would indicate clock, calendar, coordinate, or solar timing; missing contrast would suggest an unapplied multiplier, changing forcing, or insufficient pre-edge history. These alternatives are not supported by the complete PASS result.

\section{Scope and limitations}

This synthetic case covers one date, location, fuel, factor, and grid. It does not validate factor magnitude, other seasons or coordinate systems, a burn interval wrapping midnight, heterogeneity, wind, or crown-fire, spotting, suppression, and interpolation interactions. The estimator assumes circular growth, discards post-edge behavior, samples every \SI{30}{minute}, and infers rather than records the internal factor. This is implementation verification, not validation of real-fire nocturnal behavior.
